\chapter{ML and Evolutionary Samplers}
\label{ch:ml-evolutionary}

\newthought{Two samplers steer themselves} toward the interesting parts of a space rather than
covering it uniformly: \sampler{DNN} learns where to look with a neural-network surrogate, and
\sampler{Diver} evolves a population toward the best fit.

\section{DNN}
\label{sec:s-dnn}

\paragraph{How it works} \sampler{DNN} is an iterative, machine-learning-guided sampler. It starts
with an initial batch of points, then repeats a loop: train a neural network on the points seen so
far to predict which regions look promising, draw new candidates where the surrogate is
optimistic, evaluate them with the real workflow, and retrain. Over several iterations the
sampling concentrates where it matters, so expensive evaluations are not wasted on barren regions.

\paragraph{When to use} Expensive likelihoods where a surrogate can focus the search---especially
when good regions are sparse and you want to find them with fewer full evaluations.

\paragraph{Format} The \yamlkey{Bounds} block configures both the loop and the network
(Table~\ref{tab:dnn-keys}).

\begin{table}[ht]
  \footnotesize
  \begin{tabular}{@{}ll@{}}
    \toprule
    \textbf{Key} & \textbf{Meaning} \\
    \midrule
    \yamlkey{Ninit}         & number of initial points before the first training \\
    \yamlkey{Niters}        & number of guided iterations \\
    \yamlkey{Hidden\_layers} & network shape, e.g. \code{[64, 64]} \\
    \yamlkey{Batch\_size}   & training batch size \\
    \yamlkey{Nepoch}        & training epochs per iteration \\
    \yamlkey{Learning\_rate} & optimiser learning rate \\
    \yamlkey{Outputs}       & which observables the network targets \\
    \bottomrule
  \end{tabular}
  \caption{\sampler{DNN} \yamlkey{Bounds} keys.}
  \label{tab:dnn-keys}
\end{table}

\begin{yamlwin}
Sampling:
  Method: "DNN"
  Variables: [ ... ]
  Bounds:
    Ninit: 2000
    Niters: 10
    Hidden_layers: [64, 64]
    Batch_size: 128
    Nepoch: 50
    Learning_rate: 0.001
    Outputs: ["LogL"]
  LogLikelihood:
    - {name: "LogL", expression: "..."}
\end{yamlwin}

\section{Diver}
\label{sec:s-diver}

\paragraph{How it works} \sampler{Diver} uses \emph{differential evolution} (\textsc{de}) to
maximise the likelihood. It keeps a population of \yamlkey{NP} candidate points and, each
generation, builds a trial for every member by adding a scaled difference of other members
($F\cdot(x_a - x_b)$, the \yamlkey{F} mutation factor) and mixing coordinates with crossover
probability \yamlkey{Cr}; a trial replaces its parent only if it scores better. The self-adaptive
variants \code{jDE} and \code{lambdajDE} tune \yamlkey{F} and \yamlkey{Cr} automatically as the
run proceeds. This makes \sampler{Diver} robust on rugged, multi-modal landscapes.

\paragraph{When to use} Global optimisation and best-fit discovery---finding the highest-likelihood
point(s) in a difficult space---rather than mapping a posterior.

\paragraph{Format} Settings go in a \yamlkey{run} block (the alias \yamlkey{Bounds} works too);
\yamlkey{MaxWorker} is optional. The most common keys (Table~\ref{tab:diver-keys}):

\begin{table}[ht]
  \footnotesize
  \begin{tabular}{@{}ll@{}}
    \toprule
    \textbf{Key} & \textbf{Meaning} \\
    \midrule
    \yamlkey{mode}     & strategy: \code{current}, \code{jDE}, or \code{lambdajDE} \\
    \yamlkey{NP}       & population size \\
    \yamlkey{maxgen}   & maximum generations \\
    \yamlkey{maxciv}   & number of ``civilisations'' (independent restarts) \\
    \yamlkey{F}, \yamlkey{Cr} & mutation factor and crossover probability \\
    \yamlkey{convthresh}, \yamlkey{convsteps} & convergence threshold and window \\
    \yamlkey{bndry}    & boundary handling (e.g. \code{reflection}) \\
    \yamlkey{seed}     & random seed for reproducibility \\
    \bottomrule
  \end{tabular}
  \caption{Common \sampler{Diver} \yamlkey{run} keys.}
  \label{tab:diver-keys}
\end{table}

\begin{yamlwin}
Sampling:
  Method: "Diver"
  Variables: [ ... ]
  run:
    mode: jDE
    NP: 1000
    maxgen: 500
    maxciv: 1
    Cr: 0.9
    convthresh: 1.0e-3
    convsteps: 10
    seed: 1234
  LogLikelihood:
    - {name: "LogL", expression: "..."}
\end{yamlwin}

\begin{jtip}
Differential evolution is stochastic: run \sampler{Diver} a few times with different
\yamlkey{seed}s and confirm the best fit is repeatable before trusting it.
\end{jtip}
